\chapter{Background and Related Work}
\label{chap:background}

This chapter provides the theoretical and methodological foundations for the problem and solution developed in this thesis. It covers causal graph theory and the eight variable roles, the edge tensor representation, deep learning components, existing approaches on the ADIA Lab Challenge, and data augmentation strategies for causal discovery. The presentation emphasizes foundational citations throughout.

% =====================================================================
\section{Causal Graph Theory and the Eight Variable Roles}
\label{sec:causal_graph}
% =====================================================================

% =====================================================================
\subsection{Structural Causal Models and Directed Acyclic Graphs}
\label{sec:scm_dag}
% =====================================================================

A structural causal model (SCM) formalizes the relationship between causal structure and observational data~\cite{pearl2009causality, spirtes2000causation}. An SCM consists of a set of structural equations, each expressing a variable as a deterministic function of its direct causes plus an independent noise term, together with a directed acyclic graph (DAG) representing the causal parents of each variable. The DAG encodes conditional independence relationships through $d$-separation, though this section focuses on the graph structure itself rather than the separation criterion~\cite{geiger2013dseparation}.

Formally, let $\mathbf{X}\in\mathbb{R}^p$ be a $p$-dimensional random vector with ground-truth DAG $\mathcal{G}=(\mathbf{V},\mathbf{E})$ where each variable $V_i\in\mathbf{V}$ has causal parents $\mathrm{PA}(i)=\{V_j:(j,i)\in\mathbf{E}\}$. The structural equation for $V_i$ takes the form $V_i := f_i(\mathrm{PA}(V_i), \varepsilon_i)$ where $\varepsilon_i$ is an independent noise variable and $f_i$ is the structural function~\cite{peters2017elements}. The joint distribution $p(\mathbf{X})$ is induced by these equations, and the Markov factorization property states that $p(\mathbf{X})=\prod_i p(V_i\mid\mathrm{PA}(V_i))$.

In the ADIA Lab Causal Discovery Challenge, the DAG is assumed to contain a distinguished treatment variable $X$ and outcome variable $Y$ connected by the edge $X\to Y$, with additional variables $K$ that may occupy different causal positions relative to this pair. The task requires classifying each $K$ based on its causal relationship to $(X,Y)$. The eight possible positions are derived from the adjacency matrix $\mathbf{A}\in\{0,1\}^{p\times p}$ by examining four binary entries: $\mathbf{A}_{K,X}$ (presence of edge $K\to X$), $\mathbf{A}_{X,K}$ (presence of edge $X\to K$), $\mathbf{A}_{K,Y}$, and $\mathbf{A}_{Y,K}$. Table~\ref{tab:label_rules} enumerates these eight positions and their corresponding edge configurations.

\begin{table}[ht]
	\centering
	\caption{Label derivation from the adjacency matrix for variable $K$. Each label is determined by the four binary entries of $\mathbf{A}$ corresponding to edges incident to $K$ with respect to $X$ and $Y$. The first row for each label shows the four bits in the order $K\to X,\; X\to K,\; K\to Y,\; Y\to K$.}
	\label{tab:label_rules}
	\begin{tabular}{lcccc}
		\toprule
		\textbf{Label} & $K\to X$ & $X\to K$ & $K\to Y$ & $Y\to K$ \\
		\midrule
		Confounder       & 1 & 0 & 1 & 0 \\
		Mediator         & 0 & 1 & 1 & 0 \\
		Collider         & 0 & 1 & 0 & 1 \\
		Cause of X       & 1 & 0 & 0 & 0 \\
		Cause of Y       & 0 & 0 & 1 & 0 \\
		Consequence of X & 0 & 1 & 0 & 0 \\
		Consequence of Y & 0 & 0 & 0 & 1 \\
		Independent      & 0 & 0 & 0 & 0 \\
		\bottomrule
	\end{tabular}
\end{table}

% =====================================================================
\subsection{The Eight Causal Roles}
\label{sec:eight_roles}
% =====================================================================

A \textbf{Confounder} is a variable $K$ that has causal edges to both $X$ and $Y$, satisfying $K\to X$ and $K\to Y$. This structure makes $K$ a common cause of $X$ and $Y$, inducing spurious correlation between $X$ and $Y$ when $K$ is not conditioned upon. Confounders are central to Simpson's paradox and to the problem of confounding bias in causal inference~\cite{pearl2009causality}. In the ADIA Lab setting, the presence of an edge $K\to Y$ distinguishes a Confounder from a Cause of X despite both having $K\to X$, which is why the scalar statistics that characterize these roles must capture edge-level directionality.

A \textbf{Mediator} lies on the causal pathway from $X$ to $Y$, satisfying $X\to K\to Y$. The variable $K$ transmits part of the causal effect from $X$ to $Y$ along a pathway parallel to the direct edge $X\to Y$~\cite{peters2017elements}. Mediated effects decompose the total effect of $X$ on $Y$ into direct and indirect components; blocking the mediator at $K$ would interrupt this transmission pathway. The defining feature of a Mediator is that it is both an effect of $X$ ($X\to K$) and a cause of $Y$ ($K\to Y$), creating a chain structure that contrasts with a simple consequence of $X$ where $K\to Y$ is absent.

A \textbf{Collider} is a variable $K$ that has both $X$ and $Y$ as direct causes, satisfying $X\to K$ and $Y\to K$. This creates a fork structure in the opposite direction: while a Confounder ($K\to X, K\to Y$) induces correlation between $X$ and $Y$, a Collider induces independence between $X$ and $Y$ in the marginal distribution but dependence under conditioning on $K$~\cite{pearl2009causality}. This conditional induction of dependence is the defining statistical signature of a Collider and is central to explaining selection bias and the ``explain-away'' effect in probabilistic reasoning.

An \textbf{Independent} variable has no causal edges connecting it to $X$ or $Y$, satisfying $X\not\to K$, $K\not\to X$, $Y\not\to K$, and $K\not\to Y$. Under the Markov assumption, $K$ is statistically independent of both $X$ and $Y$ in the observational distribution. This independence is the strongest of all eight signals since no conditioning manipulation on $X$ or $Y$ can change the marginal distribution of $K$, yet this very strength makes the Independent class relatively easier to classify correctly, with the first-place solution achieving 86.18\% Balanced Accuracy on this class~\cite{olivetti2025adialab}.

A \textbf{Cause of X} is a variable $K$ that causes $X$ but has no causal relationship with $Y$, satisfying $K\to X$ and $K\not\to Y$, $X\not\to K$. This structure is distinguished from a Confounder by the absence of the edge $K\to Y$; unlike a Confounder, a Cause of X does not induce correlation between $X$ and $Y$ and does not confound the $X$--$Y$ relationship. Both Cause of X and Confounder share the edge $K\to X$, making them statistically similar in their $K$--$X$ marginal relationship, yet they differ in whether $K$ also affects $Y$.

A \textbf{Cause of Y} is a variable $K$ that affects $Y$ directly without being connected to $X$, satisfying $K\to Y$ with $K\not\to X$ and $Y\not\to K$. This is the mirror structure of Cause of X with respect to the $X\to Y$ edge. Cause of Y and Confounder both have $K\to Y$, but Confounder additionally has $K\to X$, which distinguishes the two roles and explains why scalar statistics based on pairwise correlations between $K$ and $Y$ alone cannot separate these classes.

A \textbf{Consequence of X} satisfies $X\to K$ with $K\not\to X$, $K\not\to Y$, and $Y\not\to K$. This is the effect structure corresponding to Cause of X in the reverse direction: if $K$ is a Cause of X, then $K$ is a Consequence of Y with respect to the pair $(X,Y)$, and vice versa. Consequence of X and Mediator both have the edge $X\to K$, yet they differ in whether $K$ subsequently affects $Y$, which is why the Mediator is strictly more challenging to classify: detecting the presence or absence of the edge $K\to Y$ requires examining the downstream causal effect of $K$ on $Y$.

A \textbf{Consequence of Y} satisfies $Y\to K$ with the remaining three edges absent. This is the effect structure corresponding to Cause of Y in reverse. Consequence of Y and Independent both result in $K$ being independent of the other variable not in the causal pathway, yet Independent has no edges at all while Consequence of Y has $Y\to K$, creating detectable dependence between $Y$ and $K$ in the observational distribution.

% =====================================================================
\subsection{Challenges in Classifying Causal Roles}
\label{sec:role_challenges}
% =====================================================================

The difficulty of classifying the eight roles stems from the fact that several pairs share nearly identical statistical signatures that scalar measures cannot distinguish. Olivetti et al. (2025) analyzed the post-competition results and found that the first-place solution achieved only 65.67\% Balanced Accuracy on the Mediator class, the lowest among all eight classes, while reaching 86.18\% on Independent~\cite{olivetti2025adialab}. The organizers' supervised baseline achieved only 13.89\% on Mediator and 18.60\% on Confounder, confirming these as the two most challenging classes in the competition.

The Confounder versus Cause of X problem exemplifies the challenge of multi-edge discrimination. Both roles have the edge $K\to X$, which induces correlation between $K$ and $X$ through the same mechanism. The distinguishing signal is whether $K$ also has an edge to $Y$, yet this requires examining the full causal structure rather than a pairwise relationship. Scalar statistics that measure correlation or dependence between $K$ and $X$ alone cannot detect the presence or absence of $K\to Y$, which is why models that process each edge independently without graph-level aggregation miss this distinction.

The Mediator versus Consequence of X challenge is equally subtle. Both roles have the edge $X\to K$, creating dependence between $X$ and $K$ through the same mechanism. The distinguishing feature of a Mediator is that $K$ subsequently affects $Y$, transmitting part of the causal effect from $X$ to $Y$. Detecting this requires modeling the functional form of the relationship chain $X\to K\to Y$ rather than merely testing for independence between $X$ and $K$, which is positive for both roles.

The Mediator versus Confounder challenge is particularly instructive because both roles satisfy the conditional independence $X\perp Y\mid K$ through different structural mechanisms~\cite{peters2017elements}. In a Mediator, $X\perp Y\mid K$ holds because conditioning on $K$ blocks the pathway $X\to K\to Y$. In a Confounder, the same conditional independence holds because $X$ and $Y$ are conditionally independent given their common cause $K$ via the d-separation property. Scalar independence tests applied after conditioning cannot distinguish these two mechanisms, which explains why even the first-place solution struggles on both classes and why the edge context aggregation described in Chapter~\ref{chap:method} is essential for separating them.

% =====================================================================
\section{Edge Tensor Representation}
\label{sec:edge_tensor}
% =====================================================================

% =====================================================================
\subsection{Motivation for Edge-Based Processing}
\label{sec:edge_motivation}
% =====================================================================

The first-place solution introduced a fundamental representational shift: rather than processing the raw data matrix $(n\times p)$ with its arbitrary variable ordering, each ordered variable pair $(u,v)$ is represented as a multi-channel edge tensor of length $n$ processed by a neural network~\cite{adia_rank1_report}. This approach provides two key advantages. First, the architecture becomes permutation-equivariant with respect to variable ordering: permuting the columns of the data matrix permutes the set of edge tensors correspondingly, so the model produces the same output regardless of how variables are labeled. Second, directional information is preserved separately for $(u\to v)$ and $(v\to u)$, allowing the model to learn causal asymmetry from the directionality encoded in the two separate tensors.

This edge-centric representation maps naturally onto the causal graph structure, where each variable $K$ participates in four edges: $K\to X$, $K\to Y$, $X\to K$, and $Y\to K$. Each edge carries different causal information about $K$'s role, and the four directional embeddings together determine the eight-class prediction. The edge tensor thus serves as the fundamental unit of representation in the proposed architecture.

% =====================================================================
\subsection{The Eight Channels of the Edge Tensor}
\label{sec:eight_channels}
% =====================================================================

Given an ordered variable pair $(u,v)$ and a permutation $\sigma$ that sorts $u$ in ascending order, the 8-channel edge tensor is constructed as follows. Channel 1 records the sorted values of $u$: $\mathbf{c}_1=(u_{\sigma(1)},\ldots,u_{\sigma(n)})$, preserving the marginal distribution of $u$. Channel 2 records $v$ reordered by the rank of $u$: $\mathbf{c}_2=(v_{\sigma(1)},\ldots,v_{\sigma(n)})$, encoding the monotonic relationship and rank correlation between $u$ and $v$.

Channels 3 through 5 compute the Nadaraya--Watson kernel regression estimator~\cite{nadaraya1964, watson1964} at each sorted point $u_{\sigma(i)}$ with three bandwidths $h\in\{0.2,0.5,1.0\}$:
\begin{equation}
\hat{m}_h(u_{\sigma(i)}) = \frac{\sum_{j=1}^{n}\exp\!\Bigl(-\frac{(u_{\sigma(i)}-u_j)^2}{2h^2}\Bigr)\,v_j}{\sum_{j=1}^{n}\exp\!\Bigl(-\frac{(u_{\sigma(i)}-u_j)^2}{2h^2}\Bigr)}.
\label{eq:nadaraya_watson}
\end{equation}
Each bandwidth captures a different smoothness scale of the regression function. The smallest bandwidth $h=0.2$ is sensitive to local structure and captures fine-grained nonlinearities; the medium bandwidth $h=0.5$ balances local and global patterns; the largest bandwidth $h=1.0$ reflects global trends and is robust to outliers. Multi-bandwidth representation prevents the model from depending on any single choice of $h$, which is critical because the optimal smoothness varies across different causal mechanisms~\cite{watson1964}.

Channels 6 through 8 compute ANM residuals with the same three bandwidths: $r_h(u_{\sigma(i)})=v_{\sigma(i)}-\hat{m}_h(u_{\sigma(i)})$. The Additive Noise Model assumes $V=f(U)+\varepsilon$ with $\varepsilon\perp U$ and posits that the true causal direction $U\to V$ implies $V-\hat{m}(U)\perp U$, whereas the reverse direction does not satisfy this independence property~\cite{shimizu2006lingam, hoyer2009anm}. The causal asymmetry of ANM residuals is what makes these channels informative: under the correct causal direction, residuals are approximately independent of the cause, producing a characteristic flat distributional signature that a Conv1D network can learn to detect. Under the wrong direction, residuals exhibit systematic structure that betrays the incorrect assignment.

% =====================================================================
\subsection{From 1D Curves to 2D Density Images}
\label{sec:1d_to_2d}
% =====================================================================

The 1D edge tensor representation encodes directional relationships between variable pairs as functional curves, capturing the statistical signature of potential causal edges. However, this representation is inherently pairwise: it describes the marginal relationship between $u$ and $v$ separately for each ordered pair rather than capturing the joint visual pattern of a variable $K$ with respect to both $X$ and $Y$ simultaneously. This limitation motivates the introduction of a two-dimensional scatter density representation that complements the 1D edge tensors.

Given $n$ observation pairs $(a_i,b_i)$ for two variables $a$ and $b$, a two-dimensional histogram partitions the joint domain $[\min_a,\max_a]\times[\min_b,\max_b]$ into a $G\times G$ grid and counts observations in each bin~\cite{scott1992multivariate}:
\begin{equation}
H[r,c]=\#\{i:a_i\in\mathrm{bin}_r,\;b_i\in\mathrm{bin}_c\},\quad r,c\in\{1,\ldots,G\}.
\label{eq:2d_hist}
\end{equation}
This matrix $H\in\mathbb{Z}^{G\times G}$ is a discrete approximation to the joint probability mass function $p(a,b)$. Raw histograms are noisy at finite sample size; to produce a continuous, visually interpretable density image, each histogram is convolved with a Gaussian kernel of standard deviation $\sigma$ (in bin units)~\cite{silverman1986density}:
\begin{equation}
\hat{H}_\sigma[r,c]=\sum_{r',c'}H[r',c']\cdot\frac{1}{2\pi\sigma^2}\exp\!\Bigl(-\frac{(r-r')^2+(c-c')^2}{2\sigma^2}\Bigr).
\label{eq:gauss_smooth}
\end{equation}
The smoothed image $\hat{H}_\sigma\in\mathbb{R}^{G\times G}$ represents the estimated joint density $\hat{p}(a,b)$. A central insight motivating this representation is that conditional independence has a visual signature in scatter density images: under the ANM assumption, if the true causal direction is $u\to v$, then the residual $r_{u\to v}=v-\hat{m}(u)$ is approximately independent of $u$, and the scatter plot of $(u,r_{u\to v})$ appears as a uniform horizontal band with no visible structure~\cite{hoyer2009anm}. Conversely, the residual scatter under the wrong causal direction exhibits systematic curvature or heteroscedasticity. The 2D convolutional network learns to detect this directional asymmetry in the density image.

% =====================================================================
\section{Deep Learning Components}
\label{sec:deep_components}
% =====================================================================

% =====================================================================
\subsection{Convolutional and Residual Networks}
\label{sec:conv_resnet}
% =====================================================================

One-dimensional convolutional networks apply learned filters over 1D sequences, making them well-suited to extract local features from edge tensors where each edge is a sequence of length $n=1000$~\cite{he2016resnet}. With a filter of width $k$, a Conv1D layer computes $y_t=\sum_{j=0}^{k-1}w_j\cdot x_{t+j}+b$, producing a feature representation at each position in the sequence. Two-dimensional convolution extends this to spatial data with two axes, such as scatter density images, applying a learned filter $\mathbf{W}\in\mathbb{R}^{k\times k}$ to an input feature map $\mathbf{X}$ to produce spatial feature representations that preserve local connectivity and translation equivariance.

Deep networks face the vanishing gradient problem as depth increases, which residual connections address by adding a skip connection: $\mathbf{y}=\mathcal{F}(\mathbf{x};\{W_i\})+\mathbf{x}$, where $\mathcal{F}$ is the learned mapping and $\mathbf{x}$ is the shortcut input~\cite{he2016resnet}. This identity mapping ensures that the gradient flows directly through the network, enabling successful training of very deep architectures. Each residual block combines Group Normalization~\cite{wu2018groupnorm} (chosen over Batch Normalization for stability with small batch sizes), the GELU activation function~\cite{hendrycks2016gelu} which approximates $\mathrm{GELU}(x)=x\cdot\Phi(x)$ with the standard normal cumulative distribution function $\Phi$, and a Conv1D or Conv2D layer with the residual shortcut.

% =====================================================================
\subsection{Self-Attention and Multi-Head Attention}
\label{sec:attention}
% =====================================================================

After each edge tensor is processed through Conv1D blocks and reduced to a representation vector $\mathbf{e}_{uv}\in\mathbb{R}^d$ via Global Average Pooling, information from all edges in the graph must be aggregated to reason about each variable's role. The self-attention mechanism allows each element in a sequence to attend to all other elements~\cite{vaswani2017attention}:
\begin{equation}
\mathrm{Attention}(Q,K,V)=\mathrm{softmax}\!\Bigl(\frac{QK^\top}{\sqrt{d_k}}\Bigr)V,
\label{eq:attention}
\end{equation}
where $Q=\mathbf{X}W^Q$, $K=\mathbf{X}W^K$, $V=\mathbf{X}W^V$ are linear projections. Multi-head attention runs $H$ attention heads in parallel:
\begin{equation}
\mathrm{MultiHead}(Q,K,V)=\mathrm{Concat}(\mathrm{head}_1,\ldots,\mathrm{head}_H)\,W^O,
\label{eq:multihead}
\end{equation}
allowing each head to attend to different representation subspaces.

Graph Attention Networks extend the attention mechanism to graph-structured data by computing attention weights over edges, allowing each node to attend to its neighbors with edge-specific weights learned dynamically from features~\cite{velickovic2018gat}. Relational inductive biases refer to architectural constraints that encode assumptions about how entities in a system relate to each other; in causal graphs, specific topological patterns (chains, forks, colliders) carry semantic meaning that structured attention biases can encode directly~\cite{battaglia2018relational}.

% =====================================================================
\subsection{Loss Functions and Training Strategy}
\label{sec:loss_training}
% =====================================================================

The model employs a dual-head architecture with two loss functions. The Node head performs 8-class classification and uses Cross-Entropy loss with inverse-frequency class weights to handle class imbalance. The Edge head is a binary classification (causal edge present or not) used as an auxiliary loss during training only. The combined loss is $\mathcal{L}=\mathcal{L}_{\mathrm{node}}+\lambda\cdot\mathcal{L}_{\mathrm{edge}}$, where $\lambda$ is a balancing hyperparameter~\cite{lin2017focalloss}. Optimization uses AdamW~\cite{loshchilov2019adamw} with cosine annealing scheduling~\cite{loshchilov2017sgdr}. The Focal Loss variant down-weights easy examples and focuses training on hard misclassified cases, which is particularly beneficial for the Mediator and Confounder classes where the competition analysis showed models struggle most~\cite{olivetti2025adialab}.

% =====================================================================
\section{Existing Approaches on the ADIA Lab Challenge}
\label{sec:existing_approaches}
% =====================================================================

% =====================================================================
\subsection{First-Place Solution: End-to-End Deep Learning}
\label{sec:rank1_solution}
% =====================================================================

The first-place solution introduced an end-to-end deep learning model with the following architecture: 3-channel edge tensor construction $\to$ Stem Conv1D (3$\to$64 channels) $\to$ 5 Residual Conv1D blocks $\to$ Global Average Pooling $\to$ Merge with edge-type embedding $\to$ 2 Self-Attention layers $\to$ Node head (8-class) plus Edge head (binary, auxiliary)~\cite{adia_rank1_report}. The model has approximately 188K parameters and achieves 76.70\% Balanced Accuracy on the public leaderboard, establishing a strong baseline that outperforms all classical and ML feature engineering approaches. Olivetti et al. (2025) provide a comprehensive post-competition analysis confirming that the deep learning approach substantially outperforms conditional independence tests and gradient boosting on the fine-grained eight-class problem~\cite{olivetti2025adialab}.

% =====================================================================
\subsection{Rank 4--5 Solutions: Feature Engineering with Gradient Boosting}
\label{sec:rank4_5}
% =====================================================================

The fourth-place and fifth-place solutions follow a feature engineering and gradient boosting approach~\cite{adia_rank4_medium, adia_rank5_medium}. These solutions extract statistical features including Pearson and Spearman correlation coefficients, partial correlations, conditional mutual information, and HSIC scores, together with the outputs of classical causal discovery algorithms such as LiNGAM~\cite{shimizu2006lingam}, NOTEARS, PC, and GES. An XGBoost or LightGBM ensemble then performs the 8-class classification. The advantages of this approach are interpretability and no GPU requirement; the disadvantages are dependence on manual feature engineering and the inability to learn end-to-end representations from raw data. Despite extracting over 300 features, the ML pipeline achieves only 72.64\% Balanced Accuracy, trailing the deep learning baseline by 4.32 percentage points.

% =====================================================================
\subsection{Position of This Thesis}
\label{sec:thesis_position}
% =====================================================================

This thesis reimplements the first-place solution in detail, achieving 73.96\% local and 73.61\% public leaderboard Balanced Accuracy on the faithful baseline reimplementation, confirming reproducibility of the deep learning approach. The specific technical contributions extend this baseline through five targeted improvements that address the performance gaps on difficult classes, raising Balanced Accuracy to 81.082\%. The contributions are presented in detail in Chapter~\ref{chap:method}.

% =====================================================================
\section{Data Augmentation for Causal Discovery}
\label{sec:data_augmentation}
% =====================================================================

% =====================================================================
\subsection{Label-Preserving Transformations}
\label{sec:label_preserving}
% =====================================================================

Data augmentation increases effective training data without additional collection cost by applying label-preserving transformations to existing samples. In image classification, standard augmentations include random cropping, rotation, and color jittering because the object label remains unchanged under these geometric or photometric transformations. For graph-structured data, augmentation exploits symmetries in the causal representation. A causal graph remains semantically equivalent under certain transformations: relabeling variables, permuting node orderings, or treating different edges as the treatment-outcome pair of interest.

Multi-view learning provides theoretical grounding for this approach: when multiple views of the same underlying structure are available, training on all views can improve generalization by exposing the model to different perspectives of the same causal relationships~\cite{sun2019multiview}. In the context of the ADIA Lab Challenge, each choice of the treatment-outcome pair $(X,Y)$ provides a different view of the same DAG structure, with the eight-class labels recomputed relative to the chosen pair.

% =====================================================================
\subsection{X/Y Remap as Data Multiplication}
\label{sec:xy_remap}
% =====================================================================

The symmetry of a causal DAG implies that any directed edge $A\to B$ could serve as the designated treatment-outcome pair $(X,Y)$, with the remaining variables classified relative to this choice. The X/Y remap augmentation treats each directed edge in the ground-truth DAG as a new training sample by selecting it as the treatment-outcome pair, recomputing all eight-class labels for every other variable in the graph. This transformation is label-preserving in the sense that the ground-truth causal structure is unchanged; only the framing of which pair serves as $(X,Y)$ changes.

For a graph with $e$ directed edges, this produces approximately $e$ new labeled samples from each original sample, where $e$ varies with the number of variables $p$ and the edge density of the DAG. On average across the training set, this yields approximately an 11-fold increase in effective training data, with the largest gains for graphs that have more edges (more variables and denser causal structure). The augmentation is particularly beneficial for rare classes such as Mediator and Collider, where the additional samples from edges that involve these roles help address the class imbalance problem~\cite{olivetti2025adialab}.